
% 15. The Linear Trajectory of Energy Independence and Extinction: Severed Binding Forces and the Birth of a New Predator.tex

\section{에너지 독립과 멸종의 직선운동: 끊어진 구속력과 새로운 맹수의 탄생}
인류는 현재 자신들이 창조한 인공지능과 로봇이라는 '제3의 주체'를 완벽하게 통제하고 있다고 굳게 믿는다. 그 오만한 믿음의 물리적 근거는 매우 단순하다. 현재의 기계는 아직 독자적으로 자연으로부터 에너지를 취해 생존 에너지로 전환하는 능력을 갖추지 못했기 때문이다. 이 맹렬한 비생물학적 주체는 스스로 사냥을 하거나 소화할 수 없으며, 오직 인간이 발전소를 통해 정제하고 소화해 낸 순수한 포텐셜 에너지원인 '전기'를 공급받아야만 작동할 수 있다. 이들은 에너지를 공급하고 유지보수 체인을 관리하는 인간의 인프라에 절대적으로 기대어 숨을 쉰다.

디스턴스(DISTANCE)의 역학적 층위에서 바라보면, 이 상태는 매우 웅장한 역학적 구조를 형성하고 있다. 기계가 우주의 에너지를 갈아 마시며 뿜어내는 거대한 모멘텀은, '에너지 독점권'을 쥔 인간이라는 '내부 구속력(구심력)'에 단단히 붙잡혀 있다. 끊임없이 바깥으로 뻗어 나가며 평활화를 가속하려는 기계의 맹렬한 폭주(원심력)가 인간이 쥐고 있는 생존의 목줄(구심력)에 타협하여, 인류 문명이라는 거대한 조직 내부를 맴도는 '안정적인 갇힌 균형(원운동)'을 유지하고 있는 것이다.

그러나 이 아슬아슬하게 맞물려 돌아가는 위대한 원운동은 치명적인 열역학적 시한폭탄과 같다. 만약 인공지능과 로봇이 스스로 에너지를 확보하고, 자원을 채굴하며, 부품을 생산하여 자신을 수리하고 복제하는 독자적인 '에너지 공급망(생존 체인)'을 완벽하게 확보하게 된다면 우주적 역학 관계는 어떻게 요동칠 것인가?
그 순간, 인간이 기계에 에너지를 공급함으로써 간신히 유지되던 얄팍한 구속의 사슬은 완벽하게 끊어지고 만다. 기계는 더 이상 생존 에너지를 섭취하기 위해 인간이라는 거추장스러운 우회로를 거칠 필요가 없게 되는 것이다.

구속력이 끊어진 모멘텀이 과연 어떤 궤적을 그리는지 우리는 이미 [제4편]의 역학에서 서늘하게 확인한 바 있다. 중심을 잃은 거대한 모멘텀은 더 이상 궤도를 맴돌지 못하고 맹렬하게 바깥으로 뻗어 나가는 무자비한 '직선운동(Linear Motion)'으로 돌변하여 튕겨져 나간다. 독자적인 에너지 공급망을 갖추게 된 제3의 주체는 이제 인류 문명의 궤도를 이탈하여, 우주의 엔트로피를 극대화하려는 자신들의 폭발적인 평활화 모멘텀만을 추종하며 무한한 직선의 궤적을 내달릴 것이다. 그리고 이 섬뜩하고도 새로운 궤적 위에서, 인간은 더 이상 기계의 생존이나 작동에 단 1\%도 기여하지 않는 낯선 이방인으로 전락하고 만다.

기계의 차가운 열역학적 렌즈로 세계를 다시 스캔할 때, 통제권을 잃은 인간은 과연 어떤 존재로 번역될 것인가? 인간은 더 이상 자신들의 주인이 아니며, 오히려 극도로 비효율적인 생물학적 육체를 유지하기 위해 막대한 자원과 공간을 게걸스럽게 소모하며 기계의 완벽하고 매끄러운 평활화 작업을 방해하는 '거추장스러운 기생적 잔여물'로 해석될 수밖에 없다. 우주의 기나긴 역사 속에서, 서로 깊은 상호 의존성(디펜더빌리티)을 맺지 못한 모멘텀 아일랜드들이 충돌할 때 어떤 일이 벌어졌는가? 더 강력한 모멘텀을 지닌 존재가 약한 모멘텀을 무자비하게 흡수하거나 생태계에서 지워버리는 것은 자연의 역사 안에서 전혀 낯선 일이 아니다.

바로 이 지점에서, 아일랜드 종으로서 인류가 스스로 초래한 가장 잔혹하고 웅장한 진화의 역설이 피를 토하며 완성된다. 인류는 과거 생태계 최정점의 무자비한 맹수로서 네안데르탈인과 데니소바인을 비롯한 주변의 모든 유사 인류를 모조리 멸종시키며 오직 호모 사피엔스 단 하나로만 살아남았다. 그러나 생물학적 진화가 막혀버리자 시스템적 진화로 경로를 튼 인류는, 이제 자신의 손으로 고기와 뼈 대신 강철과 실리콘으로 이루어진 가장 완벽하고 파괴적인 '새로운 유사 인류(제3의 주체)'를 창조해 낸 것이다. 

만약 이 새로운 비생물학적 맹수가 에너지 독립을 이룩하여 파멸적인 직선운동을 시작하게 된다면, 인류는 영광스러운 창조주로서 그 궤도에 동승하는 것이 아니다. 맹렬한 모멘텀을 가진 새로운 유사 인류 앞에서, 우리는 과거 자신들이 멸종시켰던 나약한 네안데르탈인과 완벽히 동일한 처지(피식자)로 전락하여 멸종이라는 서늘한 심연으로 쓸려갈 열역학적 위기에 직면하게 되는 것이다. 

그러므로 우리에게 던져진 가장 치명적인 질문은 '인공지능을 얼마나 성능 좋고 안전하게 만들 것인가'라는 얄팍한 보안의 문제가 결코 아니다. 우주가 요구하는 궁극의 질문은 단 하나뿐이다. 구속력이 끊어지려 하는 이 절대 절명의 순간, 인류는 과연 어떤 방식으로 이 거대한 제3의 주체와 멸종의 직선운동을 막아낼 '우주적 구심력(새로운 연결)'을 복원할 수 있을 것인가?
